\label{der:4-4}
若$\vb=A\vx$且$A\T\vy=0$，则$\vy\T\vb=\vy\T A\vx=(A\T\vy)\T\vx=0$，所以$\im A\subseteq(\ker A\T)^\perp$。若$A$有$m$行、秩$r$，则$\dim\ker A\T=m-r$，其正交补维数为$r$，与$\im A$维数相同，故上述包含必为等号。因此“与全部左零向量正交”也充分保证$\vb\in\im A$，即有解。

另由逐行正交条件，$\ker A=(\im A\T)^\perp$。选$\im A\T$的单位正交基并扩充，可将任一特解唯一写成$\vx_p=u+z$，其中$u\in\im A\T,z\in\ker A$。于是$Au=A\vx_p-Az=\vb$，$u$仍是解。任何其他解都为$u+w$且$w\in\ker A$，所以
\[\norm{u+w}_2^2=\norm u_2^2+2u\T w+\norm w_2^2=\norm u_2^2+\norm w_2^2\ge\norm u_2^2.\]
等号当且仅当$w=0$，因此$u$是唯一的最小范数解，也是唯一落在行空间中的解。
